\makeabstract
    {
    Tracking Nanoplastic Transport: Preparation of Plastics with Metal Tags and Plasticizer Additives
    }
    {
    Joel Kim\textsuperscript{1}, Elias Buurma\textsuperscript{2}, Howard Fairbrother\textsuperscript{2}
    }
    {
    \textsuperscript{1} Amherst College, Amherst, MA\\
\textsuperscript{2} Department of Chemistry, Johns Hopkins University, Baltimore, MD
    }
    {
    This study worked to broaden the application of our lab{\textquoteright}s novel top-down production of nanoplastics embedded with metal-tags (Ex: tantalum(V) ethoxide) or plasticizers (Ex: dimethyl phthalate) by testing combinations of organic solvents, commercial polymers, and additive types. With different nanoplastic variations, we aim to expand upon the characterization of nanoplastic travel, whether it be nanoplastic particles across biological systems or embedded additives when exposed to environmental stimuli.
    }
    {
    Our nanoplastics were developed by first sonicating the desired solvent, polymer, and additive to create a fully-dissolved solution, which was then left to evaporate under controlled conditions. The resulting plastic film (coupon) was cryogenically milled to generate a heterogenous mixture of nanoplastic sizes and structures. Proper additive incorporation was tested for both the coupon and the milled particles via surface analysis: Attenuated total reflectance Fourier-transform infrared (ATR-FTIR) Spectroscopy and X-ray photoelectron spectroscopy (XPS).
    }
    {
    For successful metal-tag combinations, ATR-FTIR and XPS spectra confirmed proper additive incorporation for both coupon and milled forms, as there were no associated metal signals along our sample surfaces. This meant the complete envelopment of metal tags within our nanoplastic particles, preserving the plastic{\textquoteright}s native surface properties. For successful plasticizer combinations, ATR-FTIR spectra of the coupon and the milled forms confirmed plasticizer-additive retention, as increasing plasticizer loadings corresponded to increasing plasticizer peak intensities.
    }
    {
    This study confirmed the replicability of our top-down production of nanoplastics using new solvent-polymer-additive combinations, with promising additions being dicholoromethane (DCM) as a solvent option and tantalum(V) ethoxide as a metal tag option. Our next steps for the metal-tagged nanoplastics focus on nanoplastic travel and accumulation across mice tissue and tomato plants. Future plans for the plasticized nanoplastics will study additive leaching behavior based on nanoplastic particle size and environmental stimuli.
    }

    \newpage
    